\chapter{1958:Frederick Sanger}

\label{chap:chemistry1958}

The Nobel Prize in Chemistry for the year 1958 was awarded to Frederick Sanger.

\textbf{Motivation:}

for his work on the structure of proteins, especially that of insulin

\section{Frederick Sanger}

\subsection*{Early Life and Education}

Frederick Sanger was born on 1918-08-13 in Rendcomb, England.

\subsection*{Career and Major Research}

Sanger studied at St John's College, Cambridge, and received his doctorate from Cambridge in 1943. He worked in the department of biochemistry at Cambridge and later at the Medical Research Council Laboratory of Molecular Biology. He developed methods, using the reagent dinitrofluorobenzene, to label the terminal amino groups of protein chains, and applied them to determine the complete amino acid sequence of the hormone insulin, which he published in 1955.

\subsection*{Nobel Prize-winning Work}

The work for which Frederick Sanger was awarded the Nobel Prize in Chemistry in 1958 was described by the Nobel Committee as: "for his work on the structure of proteins, especially that of insulin".

Sanger determined that insulin consists of two polypeptide chains, connected by disulfide bridges, and established the exact order of the amino acids in both chains. This was the first time the complete primary structure of a protein had been established.

\subsection*{Significance and Impact}

The sequencing of insulin demonstrated that proteins have definite, precise amino acid sequences, and it founded the field of protein sequencing. The methods Sanger developed became the standard approach for determining the order of amino acids in proteins.

\subsection*{Later Career and Honors}

Sanger continued at the Medical Research Council Laboratory of Molecular Biology, where he later developed methods for sequencing DNA, for which he shared the 1980 Nobel Prize in Chemistry with Walter Gilbert and Paul Berg. He is one of only four people to have received two Nobel Prizes. He died on 2013-11-19 in Cambridge, England.

\subsection*{Legacy}

Sanger established protein sequencing and later invented the chain-termination method of DNA sequencing, the technique used for the first complete sequencing of the human genome. His name is carried by the Sanger Institute in his honor.

\subsection*{Selected Publications}

\begin{itemize}

\item \emph{The Free Amino Groups of Insulin}, Biochemical Journal, 1945.

\end{itemize}

\clearpage

\section*{Chapter Summary}

The 1958 prize recognized the determination of protein structure. Frederick Sanger established the complete amino acid sequence of insulin at Cambridge, the first protein sequence ever determined, and his later work on DNA sequencing earned him a second Nobel Prize in 1980.

